\clearpage\thispagestyle{otFront}\phantomsection\label{references}\pdfbookmark[0]{Sources et édition}{bm-references}\otPageTitre{Sources et édition}{Références officielles · France}\fontsize{11}{15}\selectfont
\par\medskip\textbf{1. Arrêté et calendrier}\par Arrêté du 10 avril 2025 publié au BO n° 16 du 17 avril 2025 : programme de mathématiques du cycle 3, applicable en CM1 depuis la rentrée 2025-2026.\par\href{https://www.education.gouv.fr/bo/2025/Hebdo16/MENE2504620A}{\textcolor{otGreenDark}{Consulter la source officielle}}\par
\par\medskip\textbf{2. Programme officiel de mathématiques du cycle 3}\par Rubriques « Cours moyen première année » : nombres, calcul et résolution de problèmes ; grandeurs et mesures ; espace et géométrie ; données et probabilités ; proportionnalité ; initiation à la pensée informatique.\par\href{https://www.education.gouv.fr/sites/default/files/programme-de-math-matiques-pour-le-cycle-3-439827.pdf}{\textcolor{otGreenDark}{Consulter la source officielle}}\par
\par\medskip\textbf{3. Accompagnement du programme}\par Éduscol : ressources d’accompagnement du programme de mathématiques au cycle 3, dont la résolution de problèmes en quatre phases et les schémas en barres.\par\href{https://eduscol.education.gouv.fr/5712/ressources-d-accompagnement-du-programme-de-mathematiques-au-cycle-3}{\textcolor{otGreenDark}{Consulter la source officielle}}\par
\par\bigskip Sources consultées le 22 octobre 2026.
Les cours sont rédigés pour ce livre ; les exercices, défis et bilans
sont originaux.\par\medskip
Périmètre du CM1 : nombres entiers jusqu’à 999 999 ; fractions de
dénominateur au plus 20 et fractions décimales ; nombres décimaux
jusqu’aux centièmes ; divisions euclidiennes avec un diviseur à un
chiffre ; produit d’un décimal par un entier inférieur à 10 ;
conversions raisonnées par les relations entre unités, sans tableau
de conversion ; angles comparés sans rapporteur ; symétrie sur
quadrillage avec des axes horizontaux ou verticaux ; patron du cube ;
probabilités décrites avec les mots du programme, sans quantification
au-delà de « une chance sur deux ».\par\medskip
La proportionnalité se traite par des raisonnements en langage naturel
(« trois fois plus, trois fois moins »), sans tableau de proportionnalité
ni produit en croix, conformément au programme du cours moyen.
Les défis approfondissent les outils du livre et restent facultatifs.
L’ordre des chapitres est une proposition de progression.

\clearpage
